\section{Conclusion}
\label{sec:conclusion}
Repairing garbled text in generated videos is not the same problem as editing text: the source is a near miss, it is often already right, and everything around it should survive. \method treats repair as a sequence of evidence-gated decisions---leave the clip alone, lay the text out on the old glyphs or on the tracked box, transport the old appearance or not---around a rendering step whose spelling cannot drift. This makes it more reliable than the strongest generative editor on our benchmark and confines its changes to the text, but its glyphs still look rendered. Our negative results delimit what closing that gap requires: an appearance measure that sees glyph shape and orientation, letter segmentation reliable enough to reuse the characters the generator drew correctly, and generative edits accepted only when they can be verified---the step that verified selection takes, repairing more clips than either system and breaking none, at the cost of running the generative editor.
